\documentclass[10pt]{article}

% ---------- Document Identity (update these only) ----------
\newcommand{\DocStage}{}
\newcommand{\DocTitle}{Your Road to Tier One \textbf{SOF}}
\newcommand{\ProgramName}{182-Week Civilian-to-Selection Readiness Pipeline}
\newcommand{\ProgramSubtitle}{}

% ---------- Page ----------
\usepackage[legalpaper,left=0.5in,right=0.5in,top=0.6in,bottom=0.45in]{geometry}

% ---------- Typography ----------
\usepackage[T1]{fontenc}
\usepackage[utf8]{inputenc}
\usepackage{libertinus}
\usepackage{microtype}
\usepackage{parskip}
\usepackage{ragged2e}
\usepackage{textcomp}
\AtBeginDocument{\color{black}}
\AtBeginDocument{\pagecolor{white}}
\AtBeginDocument{\justifying}

% ---------- Landscape pages ----------
\usepackage{pdflscape}

% ---------- Color ----------
\usepackage[table]{xcolor}
\definecolor{StageBlue}{HTML}{1F4E79}
\definecolor{StageBlueDark}{HTML}{163A5A}
\definecolor{LightGray}{HTML}{F2F4F7}
\definecolor{MidGray}{HTML}{D9DEE5}
\definecolor{RuleGray}{HTML}{C7CED8}
\definecolor{HeaderInk}{HTML}{000000}
\definecolor{Ink}{HTML}{000000}

% ---------- Utility ----------
\usepackage{enumitem}
\sloppy
\setlength{\emergencystretch}{2em}

% ---------- Boxes / UI ----------
\usepackage[most]{tcolorbox}

\tcbset{
  charterTitle/.style={
    enhanced,
    colback=StageBlue!4,
    colframe=StageBlue,
    boxrule=1.25pt,
    arc=3pt,
    left=16pt,right=16pt,top=14pt,bottom=14pt,
    borderline north={3.2pt}{0pt}{StageBlue},
    borderline west={4pt}{0pt}{StageBlue},
    borderline south={1.25pt}{0pt}{StageBlue},
    drop shadow={black!25!white},
  },
  charterRule/.style={
    enhanced,
    colback=LightGray,
    colframe=RuleGray,
    boxrule=0.85pt,
    arc=3pt,
    left=12pt,right=12pt,top=10pt,bottom=10pt,
    borderline west={3pt}{0pt}{StageBlue},
  },
  charterBadge/.style={
    enhanced,
    colback=StageBlue,
    coltext=white,
    colframe=StageBlueDark,
    boxrule=0.6pt,
    arc=2pt,
    left=6pt,right=6pt,top=3pt,bottom=3pt,
  }
}
\newtcbox{\Badge}{nobeforeafter,charterBadge}

% ---------- Lists ----------
\setlist[itemize]{leftmargin=1.5em, itemsep=0.25em, topsep=0.25em}
\setlist[enumerate]{leftmargin=1.8em, itemsep=0.25em, topsep=0.25em}

% ---------- TOC ----------
\usepackage{tocloft}
\setcounter{tocdepth}{3}
\setlength{\cftbeforesecskip}{3pt}
\setlength{\cftbeforesubsecskip}{2pt}
\setlength{\cftbeforesubsubsecskip}{0pt}
\renewcommand{\cftsecfont}{\bfseries}
\renewcommand{\cftsecpagefont}{\bfseries}
\renewcommand{\cftsubsecfont}{\normalfont}
\renewcommand{\cftsubsecpagefont}{\normalfont}
\renewcommand{\cftsubsubsecfont}{\normalfont}
\renewcommand{\cftsubsubsecpagefont}{\normalfont}
\setlength{\cftsecindent}{0pt}
\setlength{\cftsubsecindent}{1.25em}
\setlength{\cftsubsubsecindent}{2.8em}
\setlength{\cftsecnumwidth}{2.1em}
\setlength{\cftsubsecnumwidth}{2.8em}
\setlength{\cftsubsubsecnumwidth}{3.6em}

% Remove the built-in TOC heading so only the custom heading is shown
\makeatletter
\renewcommand{\tableofcontents}{\@starttoc{toc}}
\makeatother

% ---------- Header/Footer: page number only ----------
\usepackage{fancyhdr}
\pagestyle{fancy}
\fancyhf{}
\renewcommand{\headrulewidth}{0pt}
\renewcommand{\footrulewidth}{0pt}
\fancyfoot[C]{\thepage}

% ---------- Section formatting ----------
\usepackage{titlesec}

\titleformat{\section}
  {\Large\bfseries\color{Ink}}
  {\thesection}{0.6em}{}
  [\vspace{0.25em}\textcolor{StageBlue}{\rule{\linewidth}{0.8pt}}]

\titleformat{\subsection}
  {\large\bfseries\color{Ink}}
  {\thesubsection}{0.6em}{}

\titleformat{\subsubsection}
  {\normalsize\bfseries\color{Ink}}
  {\thesubsubsection}{0.6em}{}

\titlespacing*{\section}{0pt}{0.95em}{0.35em}
\titlespacing*{\subsection}{0pt}{0.65em}{0.25em}
\titlespacing*{\subsubsection}{0pt}{0.55em}{0.2em}

% ---------- Table engine ----------
\usepackage{tabularray}
\UseTblrLibrary{booktabs}

% ---------- tabularray: GLOBAL table treatment ----------
\SetTblrInner{
  cells = {fg=black, bg=white, valign=t},
}

% ---------- Header helpers ----------
\newcommand{\Hdr}[1]{\shortstack[c]{\strut #1\strut}}

% ---------- Lines ----------
\newcommand{\TitleLine}{\textcolor{StageBlue}{\rule{0.98\linewidth}{0.95pt}}}
\newcommand{\SubTitleLine}{\textcolor{RuleGray}{\rule{0.98\linewidth}{0.65pt}}}

% ---------- Continued on next page ----------
\newcommand{\ContinuedOnNextPageInline}{%
  \par
  \vfill
  \noindent\textcolor{RuleGray}{\rule{\linewidth}{1pt}}\vspace{-2pt}
  \noindent\hfill\textit{Continued on next page}\par\vspace{-10pt}
  \noindent\textcolor{RuleGray}{\rule{\linewidth}{1pt}}
  \clearpage
}

% ---------- PDF hyperlinks ----------
\usepackage[hidelinks]{hyperref}
\hypersetup{
  colorlinks=false,
  pdfborder={0 0 0},
  linktoc=all,
  pdftitle={\DocTitle},
  pdfauthor={\ProgramName}
}

% =========================================================
\begin{document}

\section*{7.19 — Protective and Assistive Equipment}
\phantomsection
\addcontentsline{toc}{section}{7.19 — Protective and Assistive Equipment}

\subsection*{7.19.1 Purpose \& Scope}
\phantomsection
\addcontentsline{toc}{subsection}{7.19.1 Purpose \& Scope}

7.19 covers protective and assistive equipment that reduces preventable risk: visibility items, lighting, protective eyewear, conservative friction-control aids, hand protection, and comfort-only supports. It supports safe execution under variable conditions and reduces cancellations caused by avoidable minor injuries (skin tears, blisters from handles, abrasions).

\subsection*{7.19.2 Governance and Safety Boundaries}
\phantomsection
\addcontentsline{toc}{subsection}{7.19.2 Governance and Safety Boundaries}

\begin{itemize}
\item Session-level visibility requirement is mandatory when needed:
  \begin{itemize}
  \item If a session occurs in low-light or traffic-adjacent conditions, visibility controls are operationally required for that session.
  \item If visibility controls are absent, substitute to daylight or indoor conditions; do not proceed.
  \item This is a safety gating rule; it does not override phase tier timing as a “requirement,” but it does determine whether the session is allowed under conditions.
  \end{itemize}
\item Prevention, not escalation:
  \begin{itemize}
  \item protective gear supports risk control and does not justify unsafe exposure escalation.
  \end{itemize}
\item Sleeves/braces posture:
  \begin{itemize}
  \item comfort-only, optional, never required, never framed as treatment.
  \item discontinue if it worsens pain, increases swelling, or alters mechanics.
  \end{itemize}
\end{itemize}

\subsection*{7.19.3 Tier Doctrine — Applied to Protective/Assistive Equipment}
\phantomsection
\addcontentsline{toc}{subsection}{7.19.3 Tier Doctrine — Applied to Protective/Assistive Equipment}

\begin{itemize}
\item Tier 0: household-only posture; daylight-only and non-traffic routes when visibility tools are absent; minimal protective posture.
\item Tier 1: essential visibility and primary light; minimal hand protection; whistle included at Tier 1 (chosen option) as a small safety-critical signaling tool.
\item Tier 2: redundancy and weather resistance; expanded hand protection; secondary light posture.
\item Tier 3: full redundancy kit with backup lights, spare batteries/charging posture, and specialty eye protection posture where relevant.
\end{itemize}

\subsection*{7.19.4 Selection Criteria \& Fit Checks}
\phantomsection
\addcontentsline{toc}{subsection}{7.19.4 Selection Criteria \& Fit Checks}

\begin{itemize}
\item Reflectivity coverage:
  \begin{itemize}
  \item visible from multiple angles; not hidden under layers.
  \end{itemize}
\item Light runtime and mounting stability:
  \begin{itemize}
  \item stable mount without bounce; sufficient runtime for intended session duration; weather resistance.
  \end{itemize}
\item Weather resistance:
  \begin{itemize}
  \item rain/snow reduces effectiveness; select gear that stays functional when wet.
  \end{itemize}
\item Fit checks:
  \begin{itemize}
  \item headlamp does not slip; reflective vest does not flap excessively; clips do not detach during movement.
  \end{itemize}
\item Glove selection posture:
  \begin{itemize}
  \item durable palm for carries/handling; adequate dexterity for adjustments; cross-reference 7.01/7.07 friction points.
  \end{itemize}
\item Sleeve posture:
  \begin{itemize}
  \item if used, must not restrict range or encourage overuse; comfort-only.
  \end{itemize}
\end{itemize}

\subsection*{7.19.5 Setup, Safety, and Anchoring}
\phantomsection
\addcontentsline{toc}{subsection}{7.19.5 Setup, Safety, and Anchoring}

\begin{itemize}
\item Light placement:
  \begin{itemize}
  \item forward illumination via headlamp/clip light; rear visibility via reflective system; redundancy in Tier 2+.
  \end{itemize}
\item Reflective placement:
  \begin{itemize}
  \item reflectives on moving limbs improve detectability; ensure not covered by jackets (coordinate with 7.14 layers).
  \end{itemize}
\item Pre-session function check routine:
  \begin{itemize}
  \item verify lights turn on; verify battery/charge; verify mounting stability; verify whistle attachment.
  \end{itemize}
\item Route/traffic posture:
  \begin{itemize}
  \item if traffic-adjacent or low-light risk cannot be mitigated, substitute to indoor or daylight sessions.
  \end{itemize}
\end{itemize}

\subsection*{7.19.6 Maintenance, Replacement, and Storage}
\phantomsection
\addcontentsline{toc}{subsection}{7.19.6 Maintenance, Replacement, and Storage}

\begin{itemize}
\item Battery management:
  \begin{itemize}
  \item consistent charging/battery replacement posture; check before outdoor sessions.
  \end{itemize}
\item Replacement cues:
  \begin{itemize}
  \item cracked clips; strap elastic failure; reduced reflectivity; water ingress; flickering lights.
  \end{itemize}
\item Cleaning:
  \begin{itemize}
  \item remove sweat/salt residue; dry fully after wet exposure.
  \end{itemize}
\item Storage:
  \begin{itemize}
  \item keep kit staged and protected from crushing; avoid moisture damage.
  \end{itemize}
\end{itemize}

\subsection*{7.19.7 Substitution Rules — Minimal-Path Alternatives}
\phantomsection
\addcontentsline{toc}{subsection}{7.19.7 Substitution Rules — Minimal-Path Alternatives}

\begin{itemize}
\item If visibility gear is absent and conditions require it:
  \begin{itemize}
  \item substitute to daylight or indoor training environment; do not proceed in low-light/traffic-adjacent conditions.
  \end{itemize}
\item If protective hand gear is absent:
  \begin{itemize}
  \item reduce exposure that causes tearing/blistering; substitute implement selection where feasible; avoid “grip through skin failure.”
  \end{itemize}
\item Stage 2.E linkage:
  \begin{itemize}
  \item log environment hazards and substitutions in Conditions Notes.
  \item record reason posture for deviations (why the session was reslotted or moved indoors).
  \item do not claim gate-grade evidence if a test was compromised by unsafe conditions.
  \end{itemize}
\end{itemize}

\subsection*{7.19.8 Phase Integration Map — Required}
\phantomsection
\addcontentsline{toc}{subsection}{7.19.8 Phase Integration Map — Required}

\begingroup
\scriptsize
\setlength{\tabcolsep}{2pt}
\renewcommand{\arraystretch}{1.12}

\begin{longtblr}{
  width=\linewidth,
  rowhead=1,
  colspec = {
    Q[l,wd=1.05in]
    Q[l,wd=1.65in]
    X[l,3.30]
    X[l,3.80]
  },
  hlines,
  vlines,
  cells = {hsep=2pt, vsep=6pt, halign=l, valign=t},
  row{odd} = {bg=white},
  row{even} = {bg=LightGray},
  row{1} = {bg=StageBlue!15, fg=black, font=\bfseries, halign=l, valign=m},
}
\Hdr{Phase} &
\Hdr{Required Tier (from Stage 6)} &
\Hdr{What Changes / Becomes Mandatory} &
\Hdr{Notes} \\
Phase 00 &
T0 &
None &
Session-level rule still applies: if low-light/traffic-adjacent, visibility controls are operationally mandatory for that session; otherwise substitute to daylight/indoor. \\
Phase 01 &
T0 &
None &
Same session-level visibility rule applies; Tier 0 indicates no phase-level requirement yet. \\
Phase 02 &
T0 &
None &
Same session-level visibility rule applies; do not use lack of gear to justify unsafe routes. \\
Phase 03 &
T1 (End) &
Tier 1 becomes mandatory by Phase 03 (End): reflective system + primary light + whistle + minimal hand protection &
Phase 03 (End): first Tier 1 milestone in Stage 6. Session-level visibility rule remains controlling for any low-light/traffic-adjacent exposure. \\
Phase 04 &
T1 (End) &
Maintain Tier 1; hand protection becomes more operationally important as ruck/carry systems appear &
Phase 04 (E/M): protect hands during load handling; cross-reference 7.07. \\
Phase 05 &
T1 (End) &
Maintain Tier 1; longer sessions increase need for stable lighting runtime &
Phase 05 (End): ensure battery posture to prevent mid-session failure. \\
Phase 06 &
T2 (End) &
Tier 2 becomes mandatory by Phase 06 (End): secondary light redundancy + weather-resistant reflectives + expanded hand protection &
Phase 06 (End): higher consequence and winter volatility justify redundancy. \\
Phase 07 &
T2 (End) &
Maintain Tier 2; enforce stricter pre-session function checks &
Phase 07 (End): prevent avoidable failures that cause reslots. \\
Phase 08 &
T2 (End) &
Maintain Tier 2; load tolerance increases fall risk under low-light/ice conditions &
Phase 08 (End): visibility + traction posture (7.14/7.11) reduces falls. \\
Phase 09 &
T2 (End) &
Maintain Tier 2; high volume increases exposure count → more opportunities for failure &
Phase 09 (End): treat maintenance/charging as schedule-critical. \\
Phase 10 &
T2 (End) &
Maintain Tier 2; multi-domain gate weeks increase consequence of a visibility failure &
Phase 10 (End): no novelty in visibility kit in protected windows. \\
Phase 11 &
T2 (End) &
Maintain Tier 2; recovery sensitivity increases consequence of falls/skin injury &
Phase 11 (End): pre-session checks remain mandatory. \\
Phase 12 &
T3 (End) &
Tier 3 becomes mandatory by Phase 12 (End): full redundancy kit + spares; sleeves remain optional comfort-only &
Phase 12 (End): redundancy reduces probability of compromised high-consequence windows. \\
Phase 13 &
T3 (End) &
Maintain Tier 3; consolidation emphasizes low variance and high reliability &
Phase 13 (End): maintain staged redundancy; enforce session-level visibility rule. \\
\end{longtblr}

\endgroup

7.19.8 Phase Integration Map Notes:

\begin{itemize}
\item Session-level visibility requirement remains mandatory whenever conditions demand it, regardless of phase tier.
\item If any artifact treats sleeves/braces as required or as treatment, requires Stage 8 review.
\item If any phase requires 7.19 Tier 1 earlier than Phase 03 (End), requires Stage 8 review.
\end{itemize}

\subsection*{7.19.9 Risks, Contraindications, and Red Flags}
\phantomsection
\addcontentsline{toc}{subsection}{7.19.9 Risks, Contraindications, and Red Flags}

\begin{itemize}
\item Traffic and visibility hazards:
  \begin{itemize}
  \item low-light and traffic adjacency without visibility kit increases collision risk; substitute to indoor/daylight.
  \end{itemize}
\item Falls due to poor illumination:
  \begin{itemize}
  \item inadequate forward illumination increases trip/fall risk; redundancy reduces failure probability.
  \end{itemize}
\item Weather exposure risks for electronics:
  \begin{itemize}
  \item water ingress and battery failure; protect and replace; check before use.
  \end{itemize}
\item Overreliance on sleeves/braces:
  \begin{itemize}
  \item treating sleeves as treatment or justification for escalation can worsen injury risk; comfort-only posture is enforced.
  \end{itemize}
\item Red flags:
  \begin{itemize}
  \item repeated near-misses in traffic-adjacent routes → immediate route change posture and/or indoor substitution.
  \item repeated equipment failures in low-light → immediate maintenance change; do not proceed until reliability is restored.
  \end{itemize}
\item Requires Stage 8 review triggers:
  \begin{itemize}
  \item any mismatch where phases demand protective/visibility gear earlier than Stage 6 timing,
  \item any attempt to weaken the session-level visibility requirement posture.
  \end{itemize}
\end{itemize}

\subsection*{7.19.10 Compliance Checklist — Pass/Fail}
\phantomsection
\addcontentsline{toc}{subsection}{7.19.10 Compliance Checklist — Pass/Fail}

\begin{itemize}
\item \textbf{PASS}/\textbf{FAIL}: Tier tables include Tier 0–Tier 3 using required schema, generic item types only, and \$/\$\$/\$\$\$ price bands only.
\item \textbf{PASS}/\textbf{FAIL}: Session-level visibility rule is stated and substitution posture enforced (daylight/indoor when gear absent).
\item \textbf{PASS}/\textbf{FAIL}: Whistle placement is consistent: listed in Tier 1 across tables and aligned to phase map.
\item \textbf{PASS}/\textbf{FAIL}: Sleeves/braces are comfort-only, optional, never required, and not treated as treatment.
\item \textbf{PASS}/\textbf{FAIL}: Phase Integration Map aligns to Stage 6 timing; mismatches are flagged for Stage 8 review (not repaired).
\end{itemize}

\end{document}